% VeloxQuant-MLX — submission-ready draft.
% TEMPLATE: written against a generic article preamble. To submit, swap in:
%   MLSys/NeurIPS  -> \documentclass{article} + the venue .sty (neurips_2026.sty / mlsys2026.sty)
%   ACM (EuroMLSys/MLOSS-style) -> \documentclass[sigconf]{acmart}
% The body, \section structure, tables and bib are template-agnostic.
\documentclass[11pt]{article}
\usepackage[margin=1in]{geometry}
\usepackage{booktabs}
\usepackage{graphicx}
\usepackage{amsmath,amssymb}
\usepackage{url}
\usepackage{xcolor}
\usepackage{caption}
\usepackage[numbers]{natbib}

\title{VeloxQuant-MLX: A Metal-Accelerated KV-Cache Quantization Suite\\
for LLM Inference on Apple Silicon Unified Memory}
\author{Rajveer Rathod\\\texttt{rathodrajveer1311@gmail.com}}
\date{2026}

\begin{document}
\maketitle

\begin{abstract}
The key--value (KV) cache, not the model weights, is the binding memory
constraint for long-context large language model (LLM) inference on Apple
Silicon, where CPU, GPU, and Neural Engine share a single unified memory
pool. Weight quantization---the dominant focus of the local-inference
ecosystem---does nothing to bound cache growth, which scales linearly with
context length. We present \textbf{VeloxQuant-MLX}, an open-source library
that ports a suite of KV-cache quantization algorithms (TurboQuant, residual
vector quantization, VecInfer product VQ, RaBitQ, CommVQ, QJL, PolarQuant,
RateQuant mixed-precision allocation, and a new SpectralQuant variant) to
Apple's MLX framework, with hand-written, runtime-compiled Metal compute
kernels on the hot path. Our contribution is primarily a \emph{systems} one:
a unified, Metal-accelerated cache-quantization layer that plugs into
\texttt{mlx\_lm} in three lines with no change to \texttt{mlx\_lm.generate}.
On a verified single-shape microbenchmark our Metal product-VQ quantize
kernel is $6.9$--$14.7\times$ faster than the pure-MLX path
($13.2\times$ at sequence length $2048$) and reduces peak memory by
$98.4\%$ ($729\,\mathrm{MB}\!\to\!12\,\mathrm{MB}$) at the shape that
otherwise out-of-memories. VecInfer-1bit attains $16\times$ key-cache
compression on Qwen2.5-7B while \emph{matching} fp16 decode throughput, and
RaBitQ with 4-bit values attains $5.95\times$ full-KV compression on
Falcon3-7B. We report these results with explicit conditions, and we are
deliberate about negative results: the RaBitQ approximate-nearest-neighbor
\emph{search} path achieves recall@10 $\le 0.4$ and is not usable for
retrieval, and the suite is evaluated with reconstruction-fidelity and
coherence proxies rather than downstream task accuracy. We position the work
as an engineering and reproducibility artifact for the Apple-Silicon
inference community.
\end{abstract}

\section{Introduction}
A transformer caches, per layer, the key and value projections of every past
token so that attention need not recompute them. This KV cache grows linearly
with sequence length:
\begin{equation}
\mathrm{mem} = L \cdot H_{kv} \cdot d_{h} \cdot S \cdot 2_{(K,V)} \cdot 2_{\text{bytes}},
\label{eq:kvmem}
\end{equation}
where $L$ is layer count, $H_{kv}$ the number of KV heads, $d_h$ the head
dimension, and $S$ the sequence length. For Llama-3.1-8B at $S{=}8\mathrm{k}$
this is $4.2\,\mathrm{GB}$ and at $32\mathrm{k}$ it is $16.8\,\mathrm{GB}$
(fp16).

On discrete-GPU systems the cache lives in dedicated VRAM. \textbf{Apple
Silicon does not work this way}: M-series chips use \emph{unified memory},
so model weights, the KV cache, the OS, and every other process contend for
the same $16$ or $24\,\mathrm{GB}$. A 4-bit 7B model plus OS and developer
tools can consume $11$--$12\,\mathrm{GB}$ before the first token, leaving
only a few gigabytes for the cache. This makes KV-cache size, not weight
size, the practical ceiling on context length for Mac users.

The local-inference ecosystem has optimized almost everything
\emph{around} the cache. \texttt{llama.cpp}, Ollama, and LM Studio store the
KV cache at fp16 by default; experimental scalar cache quantization exists
but is not Metal-tuned and does not adapt to key distributions. Apple's own
\texttt{mlx\_lm} ships a full-fp16 cache. \textbf{Weight quantization---GGUF,
GPTQ, AWQ---is a one-time, offline operation and cannot help here.} Cache
quantization is fundamentally different: it runs at inference time, on every
token, for tensors that differ every generation.

\paragraph{Contributions.} This paper is honest about its altitude: most of
the underlying algorithms are published elsewhere and we re-implement them.
Our contributions are:
\begin{itemize}
  \item \textbf{A unified Apple-Silicon KV-cache quantization suite} exposing
  nine strategies behind one \texttt{mlx\_lm}-compatible cache interface,
  selectable by a config string, requiring three lines to enable.
  \item \textbf{Hand-written Metal compute kernels} (JIT-compiled via
  \texttt{mx.fast.metal\_kernel}) for the product-VQ hot path. The quantize
  kernel is verified at $13.2\times$ speedup and $98.4\%$ peak-memory
  reduction at the OOM-trigger shape (Section~\ref{sec:eval}).
  \item \textbf{SpectralQuant}, a KV codec that rotates keys into their SVD
  eigenbasis and applies separate codebooks to high- and low-variance
  dimensions---the most plausibly novel algorithm in the suite.
  \item \textbf{A 12-model engineering sweep} on production 4-bit models, with
  reproducible scripts, per-model result files, and figures.
  \item \textbf{Disclosed negative results}: a non-functional RaBitQ search
  path and an honest Qwen2.5-32B out-of-memory finding.
\end{itemize}
We make no claim to inventing the quantizers themselves; the contribution is
integration, Metal acceleration, and a unified Apple-Silicon evaluation.

\section{Background and Related Work}
\paragraph{Weight vs.\ cache quantization.} Weight quantization
(GPTQ~\citep{frantar2022gptq}, AWQ~\citep{lin2023awq}) compresses static
parameters offline. KV-cache quantization compresses dynamic, per-token
activations online, and the attention dot product $q\cdot k$ tolerates
approximate keys because reconstruction errors average out under the softmax.

\paragraph{KV-cache compression.} KIVI~\citep{liu2024kivi} uses per-channel
(keys) and per-token (values) 2-bit quantization. KVQuant~\citep{hooper2024kvquant}
adds non-uniform datatypes and outlier isolation. GEAR~\citep{kang2024gear}
combines low-rank and sparse residual correction.
QJL~\citep{zandieh2024qjl} replaces a sign-quantized Johnson--Lindenstrauss
sketch's normalization with an unbiased 1-bit estimator.
TurboQuant~\citep{zandieh2026turboquant} provides near-optimal-distortion
online vector quantization via random rotation plus scalar Lloyd--Max coding.
RaBitQ~\citep{gao2024rabitq} gives a theoretically bounded 1-bit code for
approximate nearest-neighbor search. CommVQ~\citep{commvq2025} makes additive
codebooks commute with RoPE. VecInfer~\citep{yao2025vecinfer} pairs a smooth
$+$ Hadamard dual transform with product VQ. RateQuant~\citep{ratequant2026}
allocates per-layer bit-widths by reverse-waterfilling.

\paragraph{Apple MLX.} MLX~\citep{mlx2023} is Apple's array framework for
unified memory, with lazy evaluation and a Metal backend.
\texttt{mx.fast.hadamard\_transform} provides an $O(d\log d)$ Metal rotation
and \texttt{mx.fast.metal\_kernel} compiles custom shaders at runtime---the
two primitives this work builds on. To our knowledge no prior library unifies
these KV-cache methods under MLX with custom Metal kernels.

\section{Method}
We summarize each family and the error analysis the repository's tests
substantiate. All methods compress \emph{keys}; some additionally quantize
values.

\subsection{TurboQuant: rotation + Lloyd--Max + QJL residual}
A random orthogonal rotation $\Pi$ (realized as a randomized Hadamard
transform, $O(d\log d)$) spreads information uniformly across dimensions, so a
single fixed scalar codebook calibrated for the resulting unit-norm
distribution suffices---no per-block scales are stored. Each key is
normalized to unit norm (the codebook is calibrated for unit norm; the scalar
norm is stored in fp16), rotated, and scalar-quantized with a Lloyd--Max
codebook using $b{-}1$ bits. A QJL 1-bit sign sketch of the residual yields
an unbiased correction to each attention score.
The MSE distortion obeys the TurboQuant bound
\begin{equation}
D_{\mathrm{mse}}(b)\ \le\ \frac{\sqrt{3\pi}}{2}\,4^{-b},
\label{eq:tqbound}
\end{equation}
which our integration test \texttt{test\_distortion\_bounds.py} verifies for
$b\in\{1,2,3,4\}$ (finite-level scalar Lloyd--Max overshoots the asymptotic
optimum by $\le 2.5\times$ at low bit-rate, as documented in the test).

\subsection{Residual vector quantization (RVQ) with analytical codebooks}
Two residual passes at 1 bit each, with codebooks precomputed analytically
from Gaussian (first pass) and Laplacian (residual) distribution theory,
require \emph{zero} calibration. On synthetic $d{=}128,\,b{=}2$ keys the
second pass lifts cosine similarity from $0.69$ to $0.98$ and SNR from
$-0.5\,\mathrm{dB}$ to $13.2\,\mathrm{dB}$ (OPTIMIZATION\_FINDINGS).

\subsection{VecInfer: smooth $+$ Hadamard dual transform, product VQ}
Per-(head, channel) smooth factors $\lambda_i=\sqrt{\max|K_i|}$ suppress
outlier channels; an orthonormal Walsh--Hadamard rotation then enables a
product VQ that preserves $q K^\top$ (the dual transform applies the inverse
to queries). Keys are split into sub-vectors, each coded against a trained
$256$-entry codebook (Lloyd k-means). This is the path our Metal kernel
accelerates (Section~\ref{sec:system}).

\subsection{SpectralQuant: SVD-basis signal/noise split}
Keys concentrate the large majority of their variance in a small fraction of
dimensions. SpectralQuant rotates keys into their eigenvector basis (SVD of a
calibration batch) and applies a high-resolution codebook to the
high-variance ``signal'' dimensions and a coarse codebook to the low-variance
``noise'' dimensions. The participation-ratio analysis and ablation are in
\texttt{spectral/} and pass their unit tests; over long contexts, aligning the
codec to the variance structure reduces error accumulation.

\subsection{RaBitQ: randomized Hadamard + sign packing + IVF}
Index build applies a randomized Hadamard rotation $H\,\mathrm{diag}(D)$,
computes per-key residuals against IVF centroids, packs $\mathrm{sign}$ into
$d/8$ bytes, and stores two scalar correction terms. Search probes
$n_{\mathrm{probe}}$ centroids, scores by a Metal packed-Hamming kernel, and
re-ranks the top-$M$ with exact fp16 dot products. We pair 1-bit keys with
4-bit MSE-quantized values for full-KV compression. \textbf{We report below
that the search/recall path does not work in this implementation}; the
\emph{memory} story does.

\subsection{CommVQ: RoPE-commutative additive VQ}
Standard product VQ fails under RoPE because
$\mathrm{quantize}(\mathrm{rotate}(x))\neq\mathrm{rotate}(\mathrm{quantize}(x))$.
CommVQ constrains each $2\times2$ block of every centroid in a RoPE plane to
the commuting form $\bigl[\begin{smallmatrix}a&-b\\b&a\end{smallmatrix}\bigr]$,
projecting onto this set after each EM M-step. Our tests verify the
constraint and reconstruction.

\subsection{RateQuant: reverse-waterfilling bit allocation}
Given per-layer sensitivities from a one-pass key-norm probe and a fractional
target $\bar b$, the closed-form reverse-waterfilling allocator (RateQuant
Theorem~2) returns integer per-layer bit-widths whose mean equals $\bar b$,
assigning more bits to sensitive layers. We fit the distortion curve
$D(b)=\alpha\beta^{-b}$ with $\beta\!\approx\!3.5$ for TurboQuant-family
codecs.

\section{System and Implementation}
\label{sec:system}
\paragraph{Metal kernels.} The pure-MLX product-VQ argmin materializes a
$[\text{chunk},n_{\mathrm{centroids}},\text{sub\_dim}]$ difference tensor,
which OOMs at large shapes. Our \texttt{vecinfer\_quantize\_metal} shader
accumulates squared distance in thread-local registers and never
materializes that tensor, giving the speedup and memory reduction in
Section~\ref{sec:eval}. Kernels are JIT-compiled on first use via
\texttt{mx.fast.metal\_kernel}; a \texttt{metal\_available()} probe and a
three-state \texttt{use\_metal\_kernels} flag (\texttt{None} auto /
\texttt{True} require / \texttt{False} force-pure) preserve a debuggable
pure-MLX path. We are explicit that the companion \emph{dequant} kernel is at
or below MLX \texttt{mx.take} parity; the win is the quantize path.

\paragraph{\texttt{mlx\_lm} integration.} Enabling compression is three lines
and leaves \texttt{mlx\_lm.generate} unchanged:
\begin{verbatim}
config = KVCacheConfig(method="turboquant_rvq", bit_width_inlier=1, seed=42)
caches = KVCacheBuilder.for_model(model, config)
model.make_cache = lambda *_a, **_k: caches
\end{verbatim}
Each cache wrapper quantizes then immediately dequantizes inside
\texttt{update\_and\_fetch}, so downstream scaled-dot-product attention sees
standard fp16 tensors. Wrappers track \texttt{compressed\_key\_bytes},
\texttt{fp16\_key\_bytes}, and assigned bits for byte-accurate accounting.

\paragraph{Architecture.} A quantizer registry, a cache factory, and a
\texttt{KVCacheBuilder.for\_model} path (which consumes per-layer bit-width
lists for RateQuant) decouple algorithm from integration. Bit-packing,
sliding-window cache, and observer utilities live in dedicated modules.

\section{Evaluation}
\label{sec:eval}
All numbers below are read directly from the repository's measured result
files; conditions are stated per result. Hardware: Apple M4, $16\,\mathrm{GB}$
unified memory, MLX $\ge 0.18$, Python 3.11/3.12. Models are 4-bit
mlx-community builds. ``Quality'' here is reconstruction fidelity (cosine
similarity, SNR) and a tokens-generated coherence proxy; we did not measure
downstream task accuracy (Section~\ref{sec:limits}).

\subsection{Metal quantize kernel}
Table~\ref{tab:metal} reports \texttt{vecinfer\_quantize\_metal} vs the
pure-MLX path ($n_{\mathrm{centroids}}{=}256$, $\text{sub\_dim}{=}8$). Speedup
grows with sequence length, from $6.9\times$ ($S{=}128$) to $14.7\times$
($S{=}8192$); the headline $13.2\times$ is the $S{=}2048$ point. At the
OOM-trigger shape ($H{=}4,S{=}4096,D{=}256,\text{sub\_dim}{=}4$) peak memory
falls from $729.3\,\mathrm{MB}$ to $12.0\,\mathrm{MB}$ ($98.4\%$).

\begin{table}[t]
\centering\small
\caption{Metal product-VQ quantize kernel vs pure-MLX, $D{=}128$.
Source: \texttt{figures/metal/results.json}.}
\label{tab:metal}
\begin{tabular}{rrrr}
\toprule
$S$ & pure (ms) & metal (ms) & speedup \\
\midrule
128  & 3.64   & 0.53  & 6.9$\times$ \\
512  & 13.47  & 1.26  & 10.7$\times$ \\
2048 & 55.10  & 4.18  & 13.2$\times$ \\
8192 & 228.56 & 15.57 & 14.7$\times$ \\
\bottomrule
\end{tabular}
\end{table}

\subsection{Compression and throughput}
Table~\ref{tab:vecinfer} reports the Qwen2.5-7B sweep. VecInfer-1bit reaches
$16\times$ key compression while matching fp16 throughput ($21.5$ vs
$21.0\,\mathrm{tok/s}$, $121$ vs $120$ tokens), attributable to the model's
strong GQA ratio (4 KV / 28 query heads). RVQ-1bit gives $7.5\times$ at
$98.5\%$ of fp16 throughput. We caution that TQ-2bit generated only $66$
tokens---aggressive scalar configs can trigger early repetition/stop on some
models.

\begin{table}[t]
\centering\small
\caption{Qwen2.5-7B-Instruct-4bit, $d_h{=}128$, $H_{kv}{=}4$, $S\!\approx\!120$.
Source: \texttt{figures/vecinfer/Qwen2.5-7B-Instruct-4bit/results.json}.}
\label{tab:vecinfer}
\begin{tabular}{lrrr}
\toprule
Config & tok/s & key comp. & tokens \\
\midrule
fp16 baseline & 21.0 & 1.0$\times$ & 120 \\
RVQ-1bit      & 20.7 & 7.5$\times$ & 112 \\
VecInfer-2bit & 21.3 & 8.0$\times$ & 120 \\
VecInfer-1bit & \textbf{21.5} & \textbf{16.0$\times$} & 121 \\
\bottomrule
\end{tabular}
\end{table}

\subsection{Full-KV compression and context (RaBitQ)}
On Falcon3-7B ($d_h{=}256$), RaBitQ 1-bit keys give a key ratio of
$11.6\times$; paired with 4-bit MSE values the full-KV ratio is
$\mathbf{5.95\times}$, constant across $S$
(\texttt{figures/RaBitQ/falcon/results.json}). KV-only memory at $32\mathrm{k}$
tokens falls from $3.67\,\mathrm{GB}$ (fp16) to $0.616\,\mathrm{GB}$. Under a
KV-only-budget linear model this implies on the order of $10^5$ tokens at an
$8\,\mathrm{GB}$ cache budget; we present that as a \emph{model-based
extrapolation}, not a measured row, and note the README's ``103k vs 17k''
figures fold in a total-RAM (weights $+$ OS $+$ cache) accounting we did not
independently reproduce.

\subsection{Mixed precision (RateQuant)}
Sensitivity ratios drive allocation: Falcon3-7B (ratio $6.48\times$) receives
$14{\times}b{=}2 + 14{\times}b{=}1$ and retains $100\%$ of fp16 throughput at
$5.22\times$ compression; Gemma3-4B (ratio $14.39\times$) receives
$3{\times}b{=}3 + 11{\times}b{=}2 + 20{\times}b{=}1$ and retains $91\%$ at
$5.22\times$ (README tables). The allocator mechanism is unit-tested; the
``$2.7\times$ lower perplexity degradation'' figure is asserted without a
result file we could locate and is therefore not claimed here as measured.

\subsection{SpectralQuant quality}
Against TurboQuant 3-bit at matched bit-width, SpectralQuant improves key
cosine similarity from $0.8329$ to $0.9072$ on Qwen2.5-0.5B ($+7.4$pp) and
from $0.7581$ to $0.8625$ on Gemma-4-4B ($+10.4$pp), with a $\sim\!5.3\times$
compression ratio (README; figures under \texttt{figures/spectral\_quant\_*/}).
We flag that no machine-readable result file accompanies these PNG figures, so
the exact deltas are reproduced from the repository's reported tables rather
than re-derived.

\subsection{Optimization journey}
Four dispatch-overhead removals (batching heads into one kernel; Hadamard
rotation; a boundary-sum quantizer replacing broadcast-argmin, bitwise-
identical to argmin at $100.0\%$ index match; cast cleanup) yield
$1.2$--$2.6\times$ throughput with zero quality regression. Memory-bound
Mistral-7B saturates at fp16 throughput; CPU-headroom Qwen3-4B reaches $92\%$
of fp16 at RVQ-2bit (OPTIMIZATION\_FINDINGS).

\section{Discussion and Limitations}
\label{sec:limits}
\paragraph{No downstream-task evaluation.} Our quality signal is
reconstruction fidelity plus a coherence proxy. We did not run LongBench, a
WikiText perplexity sweep across configs, or needle-in-a-haystack. This is the
chief threat to validity and the main work needed before a top-tier accept.

\paragraph{RaBitQ search is non-functional.} The ANN search path achieves
recall@10 of $0.0$--$0.4$ (\texttt{figures/RaBitQ/kernel/results.json}) and its
per-step latency is far above fp16; moreover the RaBitQ fp16 baselines read
$0.000\,\mathrm{ms}$, i.e.\ were not timed, so its ``speedup'' framing is a
measurement artifact. Only RaBitQ's \emph{memory} compression is sound.

\paragraph{Unified-memory specificity.} The advantage is memory, not raw
speed: the published VecInfer/RaBitQ CUDA kernel fusions do not port to Metal,
so on Apple Silicon several methods cost throughput for compression.

\paragraph{Honest negative result.} Qwen2.5-32B (4-bit, $\sim\!17.5\,\mathrm{GB}$
weights) OOMs on a $16\,\mathrm{GB}$ M4 across all configs; the memory
watchdog killed the process before a kernel panic. Cache quantization cannot
recover headroom the \emph{weights} have already consumed.

\paragraph{Reporting discrepancies we corrected.} The README's
``$13\times$ faster hot path'' applies to the \emph{quantize} kernel only
(dequant is at MLX parity); the ``$243$ tests'' badge does not match the
$314$ tests \texttt{pytest} collects at HEAD; and SpectralQuant's headline
($5.95\times$) disagrees with its per-model table ($5.33\times$). We use the
reproducible figures throughout.

\paragraph{Breadth vs.\ depth.} Throughput is characterized on one chip (M4)
at short generation lengths; behavior across M-series tiers and very long
contexts is not established.

\section{Conclusion}
VeloxQuant-MLX is an engineering and reproducibility artifact that makes
KV-cache quantization a three-line option for Apple-Silicon LLM inference,
backed by a Metal quantize kernel with a verified $13\times$ speedup and
$98\%$ peak-memory reduction at the shape that otherwise OOMs. We frame the
contribution honestly---integration, Metal acceleration, a 12-model sweep, and
the SpectralQuant codec---rather than as novel quantizers, and we disclose a
non-functional search path and the absence of downstream-task evaluation as
the clear next steps.

\bibliographystyle{plainnat}
\bibliography{references}
\appendix
\section{Reproducibility}
Environment: \texttt{pyproject.toml} (\texttt{mlx>=0.18, numpy>=1.26}),
\texttt{requirements.local.txt}; dev extras add \texttt{pytest, scipy, psutil}.
Tests: \texttt{python -m pytest veloxquant\_mlx/tests} ($314$ collected at
HEAD). Key reproductions:
\begin{verbatim}
# VecInfer per-model sweep (Table 2)
PYTHONPATH=. python benchmark_scripts/benchmark_vecinfer.py \
    --model mlx-community/Qwen2.5-7B-Instruct-4bit
# Metal quantize kernel proof (Table 1)
python scripts/metal_quantize_proof.py
# TurboQuant 4-config quality (Llama-3.2-3B)
python benchmark_scripts/benchmark_kv.py
# Distortion-bound check (Eq. 2)
python -m pytest veloxquant_mlx/tests/integration/test_distortion_bounds.py
\end{verbatim}
Seeds are fixed (\texttt{seed=42} in configs; per-test RNG seeded by $b$).
\end{document}
